\section*{Conclusion}
\addcontentsline{toc}{section}{\protect\numberline{}Conclusion}

This study was aimed to provide an in-depth examination of existing importance estimation techniques across various scenarios, with a particular focus on the impor-tance estimation of Transformer blocks to optimize the fine-tuning process.

Throughout the research, the concept of "importance" was defined as the identifi-cation of structures that can be removed without significantly affecting model performance. In particular, four types of structures withing \llm{} could be estimated with the impor-tance.

\begin{itemize}
\item \textbf{Separate Weights}: In \llms{} certain weights act as "systematic outliers" \\ significantly deviating from the average of their distribution and susbtantially affecting model performance. 
\item \textbf{Layers}: In \llms{} not all layers are of equal importance. Some layers exhibit a high level of redundancy, particularly within Multi-Head Attention module, while MLP layers are crucial for maintaining the overall effectiveness of the model.
\item \textbf{Transformer blocks}: In \llms{} Transformer blocks have varying levels of redundancy, especially in middle or later layers, with initial and final layers being more critical.
\item \textbf{Generated Architectures}: While designing Transformer models automatically, performance of predicted architectures could be estimated with importance heuristics, that approximates the model performance. 
\end{itemize}

After a comprehensive overview, several key challenges in the field of importance estimation were identified:
\begin{enumerate}
    \item \textbf{Lack of a Unified Methodology}: There is no general importance heuristic applicable across different scenarios and various contexts.
    \item \textbf{Quality Degradation}: Existing solutions often reduce computational costs but at the expense of model quality, which can degrade by at least 1-2\% or completely compromise the language modeling capabilities of \llms{}.
    \item \textbf{Computational Costs}: Most importance-based methods require extensive addi-tional inference on calibration data, significantly increasing time overhead.
\end{enumerate}

In the second part of this study, comparative analysis of existing approaches for evaluating the importance of Transformer blocks within LLMs during the fine-tuning process was conducted. This analysis not only assessed current techniques but also sought to optimize fine-tuning by selectively updating essential Transformer blocks. Key Transformer blocks-based importance approaches like \uidl{}, \ows{}, \sleb{}, \\ \shortllama{} were evaluated. As a baseline a cutting-edge \lora{} approach has been used. As reference methods the \lisa{} approach and the most prominent modifications of \lora{} -- \dora{}, \pissa{}, and \adalora{} -- were estimated. 

The 6 models fine-tuning was estimated with considered approaches, including \gpt{}, \llamaTwoSmall{}, \llamaTwoBig{}, \qwenSmall{}, \qwenMedium{}, and \qwenLarge{}. After the fine-tuning, in order to estimate the final quality of the model, zero-shot accuracy was measured on commonsense reasoning datasets: PIQA \cite{piqa}, HellaSwag \cite{hellaswag}, Winogrande \cite{winogrande}, ARC \cite{arc}, OpenbookQA \cite{openbookqa}, MMLU \cite{mmlu} using the lm-evaluation-harness package \cite{lm_eval}. Zero-shot PPL was measured on language modelling dataset WikiText2 \cite{wikitext}.

The results showed that:
\begin{itemize}
\item The \uidl{} method is the preferred option when quality maintenance is para-mount as this approach demonstrated robust performance improvements with fewer Transformer blocks than \lora{}. However, \uidl{} did not accelerate the process due to the required additional inference.
\item For scenarios where speed is prioritized over certain quality loss, the \sleb{} and \shortllama{} methods are effective as these approaches provide a 38\% reduction in fine-tuning time.
\item The \lisa{} approach is suitable when there are ample resources available for full parameter updates as it is able to accelerate the full fine-tuning process while maintaining the performance.
\end{itemize}

Given these findings, existing methods for estimating the importance of Transfor-mer blocks have potential for enhancing the fine-tuning process by updating only the essential units. However, there exist several limitations, that needs to be solved. First, enable the removal of non-essential units without degrading model quality. Second, efficiently pinpoint essential Transformer blocks without incurring additional time costs. Third, employ a general methodology that aligns well with overall model performance and could be applicable to various scenarios.


Therefore, future directions of this study within research project in Huawei include:
\begin{enumerate}
\item Develop a general importance estimation methodology that performs well across various scenarios and structures.
\item Create an importance estimation approach that efficiently identifies essential Transformer blocks without incurring additional time costs and without degra-ding performance.
\item Conduct a comparative analysis of other importance estimation techniques for weights, layers, and architectures.
\item Optimize the state-of-the-art \lora{} approach by automatically assigning \lora{} A,B adapters only to essential Transformer blocks, significantly reducing fine-tuning time while maintaining accuracy.
\end{enumerate}
